In \python, a procedure can call another procedure declared earlier.
A procedure can also call itself. It is then said to be
{\em recursive}. It is declared
with {\mcl def}, like any other procedure: there is nothing special to write. Recursive procedures \index{Recursive}are useful
for making drawings like those of figure \ref{fig:desrec} (page
\pageref{fig:desrec}). These recursive drawings have the peculiarity of being
found again and again, endlessly, in reduced form inside themselves (exactly
like the infinite image obtained between two mirrors facing each other). Since in practice it is
not possible to draw a line smaller than a pixel of the screen, the drawing stops
for very small drawings.
\begin{figure}[htbp]
 \includevspacePostScript[18.5cm]{DessinsRecursifs}{0.95}{0.95}
 \caption{Recursive drawings}\label{fig:desrec}
\end{figure}

\section{Simple recursion}

The structure of simple recursive drawings is understood by {\em induction} \index{Induction}on the size of the drawing.
We must first understand the {\em base\/}, that is, how to make a drawing of size 0 (or 1).
Then we must understand the {\em induction step\/}, that is,
how to build a drawing of size $n$ using identical drawings of sizes $n-1$,
$n-2$, \ldots\ strictly smaller than $n$.

Let us take the example of the recursive paving stones
of figure \ref{fig:pavec} (page \pageref{fig:pavec}).
\begin{figure}[htbp]
 \includevspacePostScript[6.8cm]{PavesRecursifs}{0.95}{0.95}
 \caption{Recursive paving stones}\label{fig:pavec}
\end{figure}
A paving stone of size 0 is reduced to a point (more precisely a pixel of the screen).
A paving stone of size $n$ consists of a paving stone of size $n - 1$
surrounded by a square whose side is $n$ times the grid size, as explained in
figure \ref{fig:indpavec} (page \pageref{fig:indpavec}).
\begin{figure}[htbp]
 \includevspacePostScript[8.3cm]{InductionPavesRecursifs}{0.95}{0.95}
 \caption{Inductive construction rule of the recursive paving stones}\label{fig:indpavec}
\end{figure}
The recursive procedure {\idt pavingStone} that draws a recursive paving stone receives as parameter the size {\idt n}
of the paving stone. If {\idt n} is zero, it draws a point. If {\idt n} is strictly positive then
it draws a square of side {\idt n} and then calls itself recursively to draw a
paving stone of size {\idt n - 1}.
\ \\
\input python/RecursivePavingStones.tex
\noindent The name {\idt pavingStone} can be used in the body of the procedure {\idt pavingStone}
itself: when \python\ executes the call {\idt pavingStone(n - 1)}, the declaration of {\idt pavingStone}
is finished and the name {\idt pavingStone} is therefore known. A recursive procedure is thus
declared and called exactly like the procedures we have seen so far.


The recursive paving stones can also be drawn {\em iteratively}, that is, using
a ``for'' or ``while'' loop. This would be more difficult with drawings like the draw sheet of a tennis
tournament shown in figure \ref{fig:tabtour} below.
\begin{figure}[htbp]
 \includevspacePostScript[4.3cm]{Tournoi1-2-4-8-16}{0.95}{0.95}%
 \caption{Inductive construction rule of the draw sheet\label{fig:tabtour}}
 \centering{of a tennis tournament}
\end{figure}

The draw sheet of a tennis tournament with only one entrant consists of one horizontal line
on which to write the name of this entrant, who is also the winner. The draw sheet of a
tennis tournament with {\idt n} en\-trants is obtained by drawing a horizontal line
on which to write the name of the winner, facing the point of a brace joining the
sub-sheets used to determine the two finalists. Each of the sub-sheets contains roughly
the same number of players, that is, {\idt n // 2} and
{\idt n - n // 2} players respectively (recall that {\idt //} is integer division). The recursive procedure that draws the draw sheet of a
tournament of {\idt n} players is given below and illustrated in the case where {\idt n} is 11.

 {
 %\vspace{\baselineskip}
 \vbox to 0pt{\vspace*{3cm}\includevspacePostScript[3cm]{Tournoi11}{0.95}{0.95}}%
 \input python/TournamentBracket.tex
 }
\section{Invariants}
\index{Invariant}When writing a recursive procedure it is important to think about the {\em in\-va\-ri\-ant\/}
of this procedure, that is, the hypotheses that the actual parameters
of the procedure and the \robot\ must satisfy for the drawing to be correct. Let us take the example of a
binary tree (see figure \ref{fig:arbin} below), a basic structure of
programming, which we have already met in the form of the tennis tournament draw sheet.
\begin{figure}[htbp]
 \includevspacePostScript[6.6cm]{ArbreBinaire}{0.95}{0.95}
 \caption{Inductive construction rule of a binary tree}\label{fig:arbin}
\end{figure}

\noindent The procedure that draws a binary tree is the following:

\input python/BinaryTree.tex

The first invariant is that the parameter {\idt l} of the procedure {\idt tree} must be
greater than or equal to 1. This is
true for the main call since {\idt l} $=$ 32. Any call
makes further calls only if {\idt l} $>$ 1. Since
{\idt l} $\geq 2$, we have
$(${\idt l // 2}$) \geq 1$ and the invariant is true when the procedure is called recursively
to draw the left subtree. Since
{\idt l} $= (2 \ast (${\idt l // 2}$))$
if {\idt l} is even and
{\idt l} $= (2 \ast (${\idt l // 2}$)) + 1$
if {\idt l} is odd, we have
{\idt l} $\geq (2 \ast (${\idt l // 2}$))$ hence
$(${\idt l} $- (${\idt l // 2}$)) \geq$ $(${\idt l // 2}$)$ so that
$(${\idt l} $- (${\idt l // 2}$)) \geq 1$
since $(${\idt l // 2}$) \geq 1$. The invariant is therefore also true when the
procedure {\idt tree} is called recursively
to draw the right subtree. Since the invariant is true for the main call
and remains true for all subsequent calls, it is always true.

In the same way one would check the invariant concerning the \robot, namely that before and after
drawing the binary tree, the \robot\ is placed at the root of the tree, facing north with its
pen down.

\section{Non-termination}
When using recursive procedures, it can happen that they call one another
endlessly, exactly as in the case of ``while'' loops whose execution does not terminate. The
simplest example is that of the procedure {\idt loopForever} below.

\label{prog:NonTerminaison}
\input python/NonTermination.tex
A call of {\idt loopForever()} immediately causes a new call of {\idt loop\-Forever()}, and so
on endlessly. At each call, the computer uses a little of its memory (called the
{\em stack}\index{Stack}) to remember what remains to be done afterwards (here the {\idt st()} that follows
the call). To avoid filling up the whole memory, \python\ limits the number of calls
in progress to about 1000. Beyond that, an error occurs ({\idt RecursionError}), which the
\robot's library reports with the message
`{\tt Error line 5 of NonTermination.py:\ endless recursion (stack full)}'
before the execution of the program stops (line 5 is the one with the call {\idt loopForever()} in the
body of the procedure). In general
the error is less obvious and occurs only for certain values of the parameters.
This limit also means that a correct recursive procedure cannot call
itself more than about 1000 times in a row: this is more than enough for our drawings,
since in practice the drawing stops well before that, at the level of the pixel.

\section{Mutual recursion}
Consider the serpentines of figure \ref{fig:serprec} (page \pageref{fig:serprec}). There are
two kinds of them. The serpentines
{\idt E$\scriptscriptstyle n$} have
a height $n$ and are drawn starting from their lower west corner, by drawing a
horizontal line from west to east and then a vertical line to move to the east end of the floor above.
The serpentines {\idt W$\scriptscriptstyle n$} have a height $n$ and are drawn
starting from their lower east corner, by drawing a horizontal line from east to west and then a
vertical line to move to the west end of the floor above.
\begin{figure}[htbp]
 \includevspacePostScript[4.4cm]{SerpentineRecursive}{0.95}{0.95}
 \caption{Inductive construction rule of a serpentine}\label{fig:serprec}
\end{figure}

\noindent The procedures {\idt serpentineE} and {\idt serpentineW} in the program below
draw a serpentine {\idt E$\scriptscriptstyle n$} and
{\idt W$\scriptscriptstyle n$} respectively, their formal parameter {\idt n} having as actual value
the height $n$ of the serpentine to be drawn.
\input python/RecursiveSerpentine.tex
\noindent Since a serpentine
{\idt E$\scriptscriptstyle 0$} or {\idt W$\scriptscriptstyle 0$} of height zero is empty,
the procedures {\idt serpentineE} and {\idt serpentineW} do nothing when the value $n$ of their
parameter {\idt n} is negative or zero.
To draw a serpentine {\idt E$\scriptscriptstyle n$} or {\idt W$\scriptscriptstyle n$} whose
height $n \geq 1$ is the actual value of their formal parameter {\idt n}, the procedures
{\idt serpentineE} and {\idt serpentineW} draw the
base and then above it a serpentine {\idt W$\scriptscriptstyle n-1$}
or {\idt E$\scriptscriptstyle n-1$} whose height is the value of {\idt n - 1}. Since
the procedure {\idt serpentineE} calls the procedure {\idt serpentineW} and the procedure
{\idt serpentineW} calls the procedure {\idt serpentineE}, these procedures are said to be
{\em mutually recursive\/}.

\index{Identifier}When using mutually recursive procedures in \python, one must
respect the rule that {\em every identifier must be declared before it is
used}. Yet the body of {\idt serpentineE} calls {\idt serpentineW}, which is declared
after it. This is allowed because \python\ looks up the meaning of the name {\idt serpentineW} only
when it executes this call, that is, during the execution of the main program
{\idt serpentineW(18)}, when both procedures have already been declared. It is therefore enough to
declare the two procedures one after the other, {\em before} the first call of either of them.
Each one can then call the other (and itself). If, on the contrary, {\idt serpentineE} were called
between the two declarations, execution would stop with the message
`{\tt NameError:\ name 'serpentineW' is not defined}'.

In general, declarations of mutually
re\-cur\-sive procedures have the following form:

\vfill

\index{Procedure}
\begin{center}
\begin{tabular}{@{}|p{12.5cm}|@{}} \hline
{\em Declaration of mutually recursive procedures:}\\ \hline \hline
\begin{minipage}[t]{12.5cm}
{\medskip

\mbox{\blc\ \ \ \cmt \#\ \small\rm Specification of the first procedure}\par
\mbox{\blc\ \ \ \mcl def\idt\ procedureName$_{1}$%
\footnote{\ A procedure name is an identifier.}%
(formalParam$_{1}^{1}$\footnote{\ A formal parameter is an identifier.},\ formalParam$_{1}^{2}$,\ \dots):}\par
\mbox{\blc\ \ \ \ \ \ \ \small\rm \ldots\ Instructions of the body of the procedure%
\footnote{\ Including calls of the procedures {\idt pro\-ce\-dure\-Name$_{1}$} and
{\idt pro\-ce\-dure\-Name$_{2}$}.}%
\footnote{\ The expressions in these instructions can use the formal parameters
{\idt formalParam$_{1}^{1}$}, {\idt formalParam$_{1}^{2}$}, \ldots\ of the procedure
{\idt pro\-ce\-dure\-Name$_{1}$}. }
 (indented) \ldots}\par

\medskip

\mbox{\blc\ \ \ \cmt \#\ \small\rm Specification of the second procedure}\par
\mbox{\blc\ \ \ \mcl def\idt\ procedureName$_{2}$%
(formalParam$_{2}^{1}$,\ formalParam$_{2}^{2}$,\ \dots):}\par
\mbox{\blc\ \ \ \ \ \ \ \small\rm \ldots\ Instructions of the body of the procedure%
\footnote{\ The expressions in these instructions can use the formal parameters
{\idt formalParam$_{2}^{1}$}, {\idt formalParam$_{2}^{2}$}, \ldots\ of the procedure
{\idt pro\-ce\-dure\-Name$_{2}$}.}
 (indented) \ldots}\par

\medskip

\mbox{\blc\ \ \ \cmt \#\ \small\rm Instructions of the program (calls of the procedures)}\par
\medskip}
\end{minipage}\\ \hline
\end{tabular}
\end{center}

\noindent In the same way, three or more mutually recursive procedures can be declared one after the other.

\needspace{6\baselineskip}\noindent \rule{137mm}{0.1mm}\par\nopagebreak %---------------------------------------------
\begin{exercice}{The exercise consists in drawing {\em fractals} \cite{mandelbrot1}
\cite{mandelbrot2}, that is,
drawings that are found again endlessly, in reduced form, inside themselves.

\tiret  Write a recursive \python\ procedure to build a Von Koch curve\/%
\footnote{\ Von Koch, Swedish mathematician, 1870--1924.}, as
shown in figure \ref{fig:segvonkoch} (page \pageref{fig:segvonkoch}). The curve $K_{0}$
is a straight line segment. The curve $K_{1}$ is obtained by dividing the segment $K_{0}$ into three equal
segments and replacing the middle one by an equilateral triangle without its base.
The curve $K_{n}$ is obtained in the same way by dividing into three each of the segments of the curve
 $K_{n-1}$. Write a program that draws a Von Koch snowflake, obtained by fragmenting in this way, outwards,
 the sides of an equilateral triangle. Do the same for a Von Koch anti-snowflake, for which the
 fragmentation is done towards the inside of the initial equilateral triangle. Use various fragmentations
 of a segment to obtain other kinds of snowflakes like those of figure \ref{fig:flocvonkoch}
 (page \pageref{fig:flocvonkoch}).

\tiret Program the drawings of the recursive patterns of figure \ref{fig:desrec}
(page \pageref{fig:desrec}).

\tiret Write a \python\ program that draws the Sierpinski\/%
\footnote{\ Waclaw Sierpinski, Polish mathematician, 1882--1969.}
squares (figure \ref{fig:carsierp}  page \pageref{fig:carsierp}),
triangles (figure \ref{fig:trisierp}  page \pageref{fig:trisierp}) and
pentagons (figure \ref{fig:pensierp}  page \pageref{fig:pensierp}).

\tiret Draw the plants of figure \ref{fig:vege}  (page \pageref{fig:vege}).

\tiret Draw the Sedgewick squares \cite{sedgewick1} of figure \ref{fig:sedgew}  (page \pageref{fig:sedgew}).

\tiret Draw the dragon curve of figure \ref{fig:dragon}  (page \pageref{fig:dragon})
according to the construction rules of
figure \ref{fig:consdragon}  (page \pageref{fig:consdragon}). The dragon curve of order $n$ is the one
seen when looking edge-on at a sheet of paper that has been folded $n$ times
onto itself and then unfolded at right angles.

\tiret Draw the Hilbert\/%
\footnote{\ David Hilbert, German mathematician, 1862--1943.}
curve of figure \ref{fig:hilbert}  (page \pageref{fig:hilbert}) according to the construction rules of
figure \ref{fig:conshilbert}  (page \pageref{fig:conshilbert}).

\tiret Draw the Peano\/%
\footnote{\ Giuseppe Peano, Italian mathematician, 1858--1932.}
 curve of figure \ref{fig:peano}  (page \pageref{fig:peano}) according to the construction rules of
figure \ref{fig:conspeano}  (page \pageref{fig:conspeano}).

\tiret Write a \python\ procedure that writes in base 2 a number passed as parameter,
noting that the binary representation $\beta(n)$ of a number $n$ is defined by:
\begin{list}{$\scriptscriptstyle\bullet$}{
}
  \item $\beta(-n) = -\beta(n)$ if $-n < 0$;
  \item $\beta(0) = 0$;
 \item $\beta(1) = 1$;
 \item $\beta(n) = \beta(n $\ {\idt //}\ $ 2) \beta(n$\ {\idt \%}\ $2)$ if $n > 1$.
\end{list}
Do the same in bases 3, 8 and 16 (using the digits {\idt A}, {\idt B}, {\idt C}, {\idt D},
{\idt E}, {\idt F} to denote the numbers 10, 11, 12, 13, 14, 15 respectively, and
``{\idt label(}{\cha "d"}{\idt ); translate(1, 0)}''
to write the digit \mbox{\cha d}).

\noindent $\Box$}\end{exercice}
\begin{figure}[htbp]
 \includevspacePostScript[5.1cm]{SegmentsVonKoch}{0.95}{0.95}
 \caption{Von Koch curve}\label{fig:segvonkoch}
\end{figure}
\begin{figure}[htbp]
 \includevspacePostScript[2.4cm]{RegleDragon}{0.95}{0.95}
 \caption{Construction rules of the dragon curve}\label{fig:consdragon}
\end{figure}
\begin{figure}[htbp]
 \includevspacePostScript[5.0cm]{RegleConsHilbert}{0.95}{0.95}
 \caption{Construction rules of the Hilbert curves}\label{fig:conshilbert}
\end{figure}
\begin{figure}[htbp]
 \includevspacePostScript[18.2cm]{FloconsVonKoch}{0.95}{0.95}
 \caption{Von Koch snowflakes}\label{fig:flocvonkoch}
\end{figure}
\begin{figure}[htbp]
 \includevspacePostScript[18.3cm]{CarreDeSierpinsky}{0.95}{0.95}
 \caption{Sierpinski squares}\label{fig:carsierp}
\end{figure}
\begin{figure}[htbp]
 \includevspacePostScript[16.0cm]{TriangleDeSierpinsky}{0.95}{0.95}
 \caption{Sierpinski triangles}\label{fig:trisierp}
\end{figure}
\begin{figure}[htbp]
 \includevspacePostScript[16.0cm]{PentagoneSierpinski}{0.95}{0.95}
 \caption{Sierpinski pentagons}\label{fig:pensierp}
\end{figure}
\begin{figure}[htbp]
 \includevspacePostScript[12.5cm]{Vegetaux}{0.95}{0.95}
 \caption{Fractal plants}\label{fig:vege}
\end{figure}
\begin{figure}[htbp]
 \includevspacePostScript[14.0cm]{Sedgewick}{0.95}{0.95}
 \caption{Sedgewick's fractal star}\label{fig:sedgew}
\end{figure}
\begin{figure}[htbp]
 \includevspacePostScript[7.2cm]{Dragon}{0.95}{0.95}
 \caption{Dragon curves of order 1, 2, \ldots, 8, 10 and 12}\label{fig:dragon}
\end{figure}
\begin{figure}[htbp]
 \includevspacePostScript[17.5cm]{Hilbert}{0.95}{0.95}
 \caption{Hilbert curves of orders 1 to 6}\label{fig:hilbert}
\end{figure}
\begin{figure}[htbp]
 \includevspacePostScript[12.5cm]{RegleConsPeano}{0.95}{0.95}
 \caption{Construction rules of the Peano curves}\label{fig:conspeano}
\end{figure}
\begin{figure}[htbp]
 \includevspacePostScript[12.2cm]{Peano}{0.95}{0.95}
 \caption{Peano curves}\label{fig:peano}
\end{figure}
\pagebreak[3]
\begin{corriges}{
\renewcommand{\carrenoir}{\hspace{-4mm}\rule{1.5mm}{1.5mm}\hspace{2.5mm}}
We give two programs, the others being in the folder {\idt book/\allowbreak 21\_recursion}:
\label{prog:SegmentDeVonKoch}
\input python/VonKochSegment.tex
\pagebreak[3]
\input python/DecimalToBinary.tex
\renewcommand{\carrenoir}{}
\noindent $\Box$}\end{corriges}
\clearpage

%%%%%%%%%%%%%%%%%%%%%%%%%%%%%%%%%%%%%%%%%%%%%%%%%%%%%%%%%%%%%%%%%%%%%%%%%%%
\noindent \framebox[\textwidth]{
\begin{minipage}{\myparboxwidth}
\setlength{\parindent}{2em}
\small \sf
\vspace{\abovedisplayskip}
\section*{Addenda}
 The languages {\sc algol60} \cite{algol60} and then {\sc lisp} \cite{lisp} were the first to
 introduce the use of recursion. The early controversy over the efficiency of recursive
 procedures is now a thing of the past, compilation techniques being well established.
 The idea of recursion was at the root of many advances in computer science, whether
 recursive data structures (such as trees),
 recursive algorithms (such as fast sorting algorithms) or, more recently,
 functional programming, of which \ocaml\ is an example (\python\ borrows some of its ideas).

The simplest way to understand a recursive
 procedure is to imagine that at each call the computer executes a new {\em copy} of the
 procedure.

 To show that this succession of calls terminates, it is enough to show that there
 exists an integer function of the values of the parameters which is po\-si\-tive and
 strictly decreases over successive calls.
 For example the procedure {\idt segment}
 of the program {\idt VonKochSegment} (page \pageref{prog:SegmentDeVonKoch}) must be
 called with a parameter {\idt n} that is positive or zero. If so, the recursive calls are of the form
 {\idt  segment(n - 1)} with {\idt n} $> 0$, so that over successive calls the actual
 parameters strictly decrease while remaining positive. This cannot go on forever.
 Consequently the procedure {\idt segment} does not loop when it is called with an
 initial parameter that is positive or zero.

\index{Function} {\em Functions} are to expressions what procedures are to instructions. In
 \python, a function declaration is written like a procedure declaration, but its body
 contains an instruction {\mcl return}\index{return} {\em e} which
 terminates the execution of the function and returns as result the value of the expression {\em e}. Procedures
 are simply functions without {\mcl return} whose result we are not interested in, which is
 the special value {\mcl None} (nothing).

\vspace{\abovedisplayskip}
\end{minipage}
}
%%%%%%%%%%%%%%%%%%%%%%%%%%%%%%%%%%%%%%%%%%%%%%%%%%%%%%%%%%%%%%%%%%%%%%%%%%%%

%%%%%%%%%%%%%%%%%%%%%%%%%%%%%%%%%%%%%%%%%%%%%%%%%%%%%%%%%%%%%%%%%%%%%%%%%%%
\noindent \framebox[\textwidth]{
\begin{minipage}{\myparboxwidth}
\setlength{\parindent}{2em}
\small \sf
\vspace{\abovedisplayskip}


 For example, a constant is a function without parameters, so
 that a declaration such as:

\smallskip
\mbox{\blc  \blk\idt  gridSize\ =\ 15}\par
\smallskip

\noindent can be replaced by a function declaration:

\smallskip
\mbox{\blc  \blk\idt  \mcl def\idt \ gridSize():\ \mcl return\idt \ 15}\par
\smallskip

\noindent which is used by writing {\idt gridSize()}, like {\idt valGsX()}.
As a second example, consider the computation of $2^{n}$ with $n \geq 0$, observing for the base
that $2^{0} = 1$ and for the induction step that $2^{n} = 2 \ast 2^{n-1}$:

\smallskip
\mbox{\blc  \blk\idt  \mcl def\idt \ twoToThePower(n):}\par
\mbox{\blc  \blk\blk\idt  \mcl if\idt \ n\ ==\ 0:}\par
\mbox{\blc  \blk\blk\blk\idt  \mcl return\idt \ 1}\par
\mbox{\blc  \blk\blk\idt  \mcl else\idt :}\par
\mbox{\blc  \blk\blk\blk\idt  \mcl return\idt \ 2\ *\ twoToThePower(n\ {-}\ 1)}\par
\smallskip

%\vspace{\abovedisplayskip}
%\end{minipage}
%}
%%%%%%%%%%%%%%%%%%%%%%%%%%%%%%%%%%%%%%%%%%%%%%%%%%%%%%%%%%%%%%%%%%%%%%%%%%%%

%%%%%%%%%%%%%%%%%%%%%%%%%%%%%%%%%%%%%%%%%%%%%%%%%%%%%%%%%%%%%%%%%%%%%%%%%%%
%\noindent \framebox[\textwidth]{
%\begin{minipage}{\myparboxwidth}
%\setlength{\parindent}{2em}
%\small \sf
%\vspace{\abovedisplayskip}
\noindent After this declaration, the expression {\idt twoToThePower(2)} is 4 and
the ex\-pres\-sion {\idt twoToThePower(twoToThePower(2))} is 16 (the program
{\idt TwoToThePower} in the folder {\idt book/\allowbreak 21\_recursion} draws the value of
{\idt twoToThePower(16)}). Here {\mcl if}\ldots{\mcl else} chooses between two
{\mcl return} instructions, hence between the two values that the function can return.

 In \python\ one can use the trigonometric functions {\idt sin} (sine),
 {\idt cos} (cosine), {\idt tan} (tangent) and {\idt atan} (arc tangent). Angles are expressed in
 radians (an angle of $d$ degrees is {\idt d / 180 * pi} radians, the constant {\idt pi}
 being $\pi \simeq 3.14159$). As for the
 functions {\idt exp} (exponential) and
 {\idt log} (natural logarithm), the
  parameter is a number (integer or rational) and the
result is of type {\idt float}. These functions and the constant {\idt pi} are provided by the
\robot's library (they come from \python's {\idt math} module). For powers, one uses
 {\idt a ** x}, which is $a^{x} = e^{x \ln a}$. One could define the tangent function oneself by:

{
\parindent \parindentprog
\medskip
\mbox{\blc  \blk\idt  \mcl def\idt \ tg(x):\ \mcl return\idt \ sin(x)\ /\ cos(x)}\par
}

\noindent When {\idt cos(x)} is {\idt 0.0}, the division causes an error and
stops the program with the message `{\tt ZeroDivisionError:\ float division by zero}'.
Because of rounding errors, this does not happen for
{\idt tg(pi / 2)}: {\idt cos(pi / 2)} is {\idt 6.123233995736766e-17} and not {\idt 0.0}, so
that {\idt tg(pi / 2)} is {\idt 1.633123935319537e+16}, a very large number instead of infinity.

\vspace{\abovedisplayskip}
\end{minipage}
}
%%%%%%%%%%%%%%%%%%%%%%%%%%%%%%%%%%%%%%%%%%%%%%%%%%%%%%%%%%%%%%%%%%%%%%%%%%%%

